%%%% 1. DOCUMENTCLASS %%%%
\documentclass[journal=tches,preprint,floatrow]{iacrtrans}

%%%% 2. PACKAGES %%%%
\usepackage{booktabs, siunitx, multirow}
\usepackage[ruled,vlined]{algorithm2e}
\usepackage{caption}
\usepackage{etoolbox}

%%%% 3. AUTHOR, INSTITUTE %%%%
\author[Adhikari, Lyubashevsky, Schwabe]{%
  Swechchha Adhikari\inst{1}\thanks{Work done while at MPI-SP, Bochum, Germany.} \and
  Vadim Lyubashevsky\inst{2} \and
  Peter Schwabe\inst{3,4}%
}
\institute{%
  Brigham Young University, Provo, USA,
  \email{adhikar6@byu.edu}
  \and
  IBM Research, Zurich, Switzerland,
  \email{VAD@zurich.ibm.com}
  \and
  Max Planck Institute for Security and Privacy, Bochum, Germany,
  \email{peter@cryptojedi.org}
  \and
  Radboud University, Nijmegen, The Netherlands%
}

%%%% 4. TITLE %%%%
\title{Evaluating ML-KEM at (much) higher security levels}

%%%% 5. SPACING TWEAKS %%%%
\geometry{hmargin=1.6cm, vmargin=0.6cm}
\makeatletter
% Kill the 2em gap above the title and shrink the title-block gaps further
\patchcmd{\@maketitle}{\vskip 2em}{\vskip 0pt}{}{}
\patchcmd{\@maketitle}{\vskip 1.5em}{\vskip 0.3em}{}{}   % title -> authors
\patchcmd{\@maketitle}{\vskip 1em}{\vskip 0.2em}{}{}      % authors -> institute
\patchcmd{\@maketitle}{\vskip 1.5em}{\vskip 0.3em}{}{}    % institute -> body
\makeatother
% Tighten space around floats and displays
\setlength{\textfloatsep}{4pt plus 1pt minus 1pt}
\setlength{\intextsep}{4pt plus 1pt minus 1pt}
\setlength{\abovedisplayskip}{2pt}
\setlength{\belowdisplayskip}{2pt}
\captionsetup{skip=2pt}

\begin{document}

\maketitle

\noindent
``Conservative'' post-quantum security could mean avoiding algebraic structure rendering the scheme vulnerable to hypothetical future attacks \cite{takeoff}, or maintaining a wide security margin against improving cryptanalysis. FrodoKEM, standardized by ISO \cite{iso}, provides the first by design, but widening its margin requires scaling parameters that already cost nearly ten times the bandwidth of NIST-standardized ML-KEM at the same security level \cite{fips203, frodo}. We show that a module scheme can buy this margin cheaply: \textsc{Kaiburr} reaches a large security margin within FrodoKEM's bandwidth while running roughly $3\times$ faster. Fixing the ML-KEM base ring, noise distribution, and
modulus, and varying only $k$, we measure the decryption-failure probability using the \textsc{Kyber} scripts~\cite{kyber-estimates}. The $2^{-128}$ bound holds at $k=7$ but fails at $k=8$, so raising $k$ further cannot buy higher security without sacrificing correctness.

Instead, one could reduce the failure probability at large $k$ by increasing the modulus $q$. However, that would mean \textsc{Kaiburr} would not be able to reuse any existing optimized routines in software or hardware. To preserve the essence of \textsc{Kyber} i.e. the arithmetic in the ring $\mathbb{Z}_{3329}[X]/(X^{256}+1)$, we propose a generalization of ML-KEM's error distribution as $f(t)$ on $\{-2,\dots,2\}$ with
$\Pr[\pm2]=2^{-t}$, $\Pr[\pm1]=\tfrac14$, $\Pr[0]=\tfrac12-2^{-(t-1)}$. A sample
is drawn from uniform bits $a_1,\dots,a_t$ as
$x=(2a_t-1)\bigl(1-a_1+2\textstyle\prod_{i=1}^{t-1}a_i\bigr)$, which is
branchless; we machine-check both CT and speculative
constant-time (SCT) in Jasmin \cite{jasmin}.

Larger $k$ raises decryption error, so we jointly raise $t$ (lower $\sigma$)
to compensate, while favoring vectorization-friendly values of $t$. $t=4$ recovers ML-KEM's sampler $\mathrm{CBD}(2)$, $t=8$ allows one-byte-per-coefficient SIMD packing, and $t=6$ gives more security per dimension and fewer PRF bytes at the cost of less regular packing. For each $t$ we take the largest $k$ meeting the failure bound and achieving better computational and bandwidth performance than FrodoKEM-640-SHAKE.

\begin{table}[H]
  \centering
  \footnotesize
  \setlength{\tabcolsep}{4pt}
  \caption{\textsc{Kaiburr} parameter sets. Classical/quantum security (Core-SVP) and failure probability computed with the Kyber scripts~\cite{kyber-estimates}; other parameters as in ML-KEM, compression removed.\protect\footnotemark}
\label{tab:kaiburr-params}
  \sisetup{group-separator={,}, group-minimum-digits=4}
  \begin{tabular}{@{}lcc c
      S[table-format=3.0] S[table-format=3.0]
      S[table-format=4.0] S[table-format=4.0]@{}}
    \toprule
    Scheme & $k$ & noise & failure & {class.} & {quant.} & {pk (B)} & {ct (B)} \\
    \midrule
    kaiburr-1792 & 7  & $f(4)$ & $2^{-133.7}$ & {479} & {434} & 2720 & 3072 \\
    kaiburr-4608 & 18 & $f(6)$ & $2^{-134.1}$ & {1281} & {1162} & 6944 & 7296 \\
    kaiburr-6144 & 24 & $f(8)$ & $2^{-139.9}$ & {1717} & {1557} & 9248 & 9600 \\
    \bottomrule
  \end{tabular}
\end{table}

\footnotetext{Removing compression allows us to obtain formally verified correctness bounds as for FrodoKEM in \cite{bound}, left as future work.}

\textsc{Kaiburr} has reference and AVX2 implementations in both C/asm and Jasmin, available \href{https://github.com/cryptojedi/kaiburr}{here}. We benchmark on two Intel generations following the methodology of \cite{epiv}: single-core with Hyper-Threading and Turbo Boost disabled, medians over \num{10001} runs with randomness from Linux \texttt{getrandom} (Table ~\ref{tab:kaiburr-cycles}). 
\begin{table}[H]
  \centering
  \footnotesize
  \setlength{\tabcolsep}{5pt}
\caption{\textsc{Kaiburr} (AVX2) vs.\ FrodoKEM-640-SHAKE (AVX2), CPU cycles;
  \textsc{Kaiburr}-6144 runs ${\sim}2.6$--$3.1\times$ faster than Frodo. Jasmin stays within 17\% of C/asm, faster on Coffee Lake and slower on Raptor Lake.}
  \label{tab:kaiburr-cycles}
  \sisetup{group-separator={,}, group-minimum-digits=4}
  \begin{tabular}{@{}lll
      S[table-format=7.0]
      S[table-format=7.0]
      S[table-format=7.0]@{}}
    \toprule
    Microarch. & Scheme & Impl. & {keypair} & {enc} & {dec} \\
    \midrule
    \multirow{7}{*}{\shortstack[l]{Coffee L.\\(i7-8700)}}
      & \multirow{2}{*}{kaiburr-1792}
        & C      & 155113  & 159177  & 171924  \\
      & & Jasmin & 144489  & 145158  & 150135  \\
      & \multirow{2}{*}{kaiburr-4608}
        & C      & 806319  & 815330  & 842327  \\
      & & Jasmin & 762472  & 770586  & 781276  \\
      & \multirow{2}{*}{kaiburr-6144}
        & C      & 1368794 & 1378658 & 1416170 \\
      & & Jasmin & 1294067 & 1303213 & 1311650 \\
    \cmidrule(l){2-6}
& FrodoKEM-640-SHAKE\footnotemark & \cite{frodo} & 3898666 & 4063179 & 3939567 \\
\midrule
    \multirow{7}{*}{\shortstack[l]{Raptor L.\\(i9-13900K)}}
      & \multirow{2}{*}{kaiburr-1792}
        & C      & 120886  & 124488  & 134556  \\
      & & Jasmin & 141676  & 143442  & 148714  \\
      & \multirow{2}{*}{kaiburr-4608}
        & C      & 647542  & 651956  & 676062  \\
      & & Jasmin & 746326  & 750672  & 763608  \\
      & \multirow{2}{*}{kaiburr-6144}
        & C      & 1095258 & 1103782 & 1133918 \\
      & & Jasmin & 1265650 & 1271618 & 1285842 \\
    \cmidrule(l){2-6}
      & FrodoKEM-640-SHAKE & \cite{frodo} & 3121248 & 3246262 & 3189934 \\
    \bottomrule
  \end{tabular}
\end{table}
\footnotetext{We modified the Frodo harness~(\url{https://github.com/microsoft/PQCrypto-LWEKE})
to take the median of \num{10001} runs for consistency.}

We also implement and report \textsc{Kaiburr} on Pavona's Asymmetric
Cryptography Coprocessor (ACC)~\cite{pavona}, which reuses ML-KEM's
arithmetic core and accelerates Keccak in hardware. We use this platform to
evaluate how scaling ML-KEM affects performance when the dominant
computational components are accelerated. Our simulation model raises
DMEM from 32 to 128~KiB for large $k$. On ACC, \textsc{Kaiburr} scales
predictably: increasing from $k=7$ to $k=24$ raises
key generation, encapsulation, and decapsulation costs from roughly $240$K
cycles to $2.5$M cycles.

This shows that scaling the security of \textsc{Kyber} is straightforward up to $k=7$ and can be achieved beyond that only by generalizing the noise distribution.

{\footnotesize
\bibliographystyle{alpha}
\bibliography{proposal}
}

\end{document}